\label{ch:nsi-results}

This chapter reports the actual results produced by the notebooks that
exercise NSI: two physics-validation notebooks
(\texttt{notebooks/physics/experimental\_results/BSM/NSI1\_test.ipynb},
\texttt{NSI2\_lma\_dark.ipynb}), one cross-method intrinsic-validation
notebook (\texttt{notebooks/validation/intrinsic/intrinsic2\_BSM\_extension\_NSI.ipynb}),
and one real-data inference notebook
(\texttt{notebooks/inference/inference2\_borexino\_nsi.ipynb}).

\begin{notebox}
\textbf{On the completeness of the reported numbers.} Several of these
notebooks' saved \code{.ipynb} files have code cells with
\code{execution\_count: null} and an empty output list, even though later
cells or their own markdown narrative reference a result -- almost
certainly because the file was saved from a session whose full re-run
output was not re-serialized for every cell. This chapter states explicitly,
for every quoted number, whether it is a value directly recorded in the
notebook's saved output (\textbf{confirmed}) or a value only implied by
later cells/markdown prose without its own recorded output
(\textbf{unconfirmed, not independently verifiable from the file as saved}).
No number is stated without this distinction. This is a bookkeeping
limitation of the notebook files as currently saved, not a defect in the
underlying \texttt{tpeanuts} computation itself, which \cref{ch:nsi-implementation}
documents and this project's automated test suite (\pkgpath{core/BSM/test/test2\_bsm\_nsi.py}
and related) independently exercises on every change.
\end{notebox}

\section{\texttt{NSI1\_test.ipynb}: NSI deformation of the solar MSW resonance}

\begin{implbox}
This notebook builds the solar propagation Hamiltonian directly via
\pyfunc{hamiltonian\_flavour} for a battery of NSI presets, diagonalises it
with \code{torch.linalg.eigh}, and integrates over the solar production
profile -- explicitly described in the notebook as "the same approach used
internally by \code{solar\_probability\_state} when \code{epsilon} is
provided." Presets exercised in its main results section: \code{sm\_no\_nsi},
\code{nsi\_diagonal\_biggio2009}, \code{nsi\_offdiag\_icecube2022},
\code{nsi\_dune\_etau}, \code{nsi\_hyperk\_etau}, and
\code{nsi\_globalfit\_esteban2018} (\cref{sec:presets-table}).
\end{implbox}

\begin{resultbox}
\textbf{Confirmed, directly printed output:}
\begin{align*}
  \text{SM closure error }(\epsilon=0\text{ vs.\ analytic}) &= 3.879\times10^{-6} \\[4pt]
  \Delta P_{ee}\text{ at }E=10\text{ MeV, relative to SM:} \\
  \code{nsi\_diagonal\_biggio2009} &: -0.0121 \\
  \code{nsi\_offdiag\_icecube2022} &: +0.0000 \\
  \code{nsi\_dune\_etau} &: -0.0201 \\
  \code{nsi\_hyperk\_etau} &: -0.0202 \\
  \code{nsi\_globalfit\_esteban2018} &: -0.0228 \\[4pt]
  \text{Diagonal } \epsilon_{ee}\text{ scan, } 30\text{ points in }[-0.10,0.90]&: \\
  \max P_{ee} \text{ in scan} &= 0.5248 \\
  \min P_{ee} \text{ in scan} &= 0.2910
\end{align*}
The SM closure error of $3.9\times10^{-6}$ compares the fully numerical
$\epsilon=0$ Hamiltonian-diagonalisation path against the closed-form
adiabatic-approximated analytic path; the notebook's own interpretation is
that this is "the intrinsic accuracy of the analytic path's mixing-angle
factorisation... not floating-point-level" -- i.e.\ the two are independent
implementations of the same physics agreeing to five significant figures,
not the same code path compared to itself.
\end{resultbox}

\begin{notebox}
The $\epsilon_{\mu\tau}$-only preset (\code{nsi\_offdiag\_icecube2022})
produces $\Delta P_{ee}=0.0000$ at 10~MeV to the precision printed: a purely
$\mu\tau$ off-diagonal coupling has no first-order effect on $\nu_e$
survival at this energy/baseline within the printed precision, consistent
with $\nu_e$ appearing only in the $ee$/$e\mu$/$e\tau$ entries of $\epsilon$,
none of which this preset touches.
\end{notebox}

The notebook produces four figure sets (\code{sn5\_fig3}--\code{sn5\_fig6}):
an SM-validation panel ($P_{ee}(E)$ overlay plus the closure-error semilog
plot above), a multi-preset $P_{ee}(E)$ comparison with a $\Delta P_{ee}$
panel, a 2-D heatmap of $P_{ee}(E,\epsilon_{ee})$ over the diagonal scan
plus 1-D slices, and an off-diagonal $\epsilon_{e\mu}$/$\epsilon_{e\tau}$
scan. The notebook's stated physical reading of these figures: diagonal
$\epsilon_{ee}$ "simply rescales the effective matter potential...shifting
the resonance energy without creating new crossings," while off-diagonal
entries "introduce new level crossings...that have no SM analogue," and
"current constraints ($|\epsilon_{e\alpha}|\lesssim0.3$--$0.5$) allow
$\mathcal{O}(5$--$15\%)$ shifts in $P_{ee}$" -- consistent with the
confirmed $\Delta P_{ee}$ values above, all in the few-percent range for
$|\epsilon|\sim0.1$--$0.3$ couplings.

\section{\texttt{NSI2\_lma\_dark.ipynb}: the LMA-Dark degeneracy}
\label{sec:results-lma-dark}

\begin{implbox}
This notebook validates the LMA-Dark degeneracy of \S\ref{sec:theory-lma-dark}
end to end: the analytic $\theta_{12}\leftrightarrow\epsilon_{ee}$ degeneracy
identity, the daytime solar $P_{ee}(E)$ curves via
\pyfunc{solar\_probability\_state}, and -- the observable that actually
breaks the degeneracy -- the Earth day--night asymmetry $A_{DN}$ via
\pyfunc{earth\_probability\_state\_numerical} (which the notebook notes
supports NSI directly through its \code{epsilon=...} argument, unlike the
daytime solar pipeline used here, which does not yet include NSI inside the
solar interior for this specific comparison).
\end{implbox}

\begin{resultbox}
\textbf{Confirmed, directly printed output} (configuration/degeneracy check):
\begin{align*}
  \theta_{12}^{\rm SM} &= 33.4100^\circ, &
  \theta_{12}^{\rm Dark} &= 56.5900^\circ \\
  \sin^2\theta_{12}^{\rm SM} &= 0.3032, &
  \cos^2\theta_{12}^{\rm Dark} &= 0.3032 \quad\text{(must match -- confirmed match)} \\
  \epsilon_{ee}^{\rm SM} &= 0.0, &
  \epsilon_{ee}^{\rm Dark} &= -2.0
\end{align*}
\textbf{Confirmed, directly printed output} (day--night asymmetry at
$E=10$~MeV, linearly interpolated on the Earth-propagation energy grid):
\begin{center}
\begin{tabular}{@{}lll@{}}
\toprule
Scenario & $A_{DN}$ & Interpretation printed by the notebook \\
\midrule
SM LMA & $+4.84\%$ & expected $>0$ \\
SM + NSI Earth & $-1.59\%$ & sign flipped \\
LMA-Dark approximation & $-2.85\%$ & LMA-Dark signature ($<0$) \\
\midrule
Super-Kamiokande (external) & $+3.3\pm1.0\%$ & favours LMA, disfavours LMA-Dark at $\sim3\sigma$ \\
\bottomrule
\end{tabular}
\end{center}
\end{resultbox}

\begin{notebox}
The confirmed \texttt{tpeanuts} SM-LMA value, $A_{DN}=+4.84\%$ at 10~MeV, is
numerically above the quoted Super-Kamiokande measurement,
$+3.3\pm1.0\%$~\parencite{Esteban2018LMADark}: both share the same sign
(qualitatively consistent, both $>0$), but the notebook itself does not
discuss the $\sim1.5$ percentage-point quantitative gap. Super-K's
$A_{DN}$ is a flux-averaged, detector-specific measurement, whereas the
$+4.84\%$ figure here is evaluated at a single energy and nadir angle
(30$^\circ$, 1~km detector depth) without folding a detector response or a
flux/exposure average, so the two are not expected to agree in detail; this
is a genuine open comparison for anyone extending this notebook, not a
resolved validation.
\end{notebox}

The notebook produces three figures: MSW eigenvalue curves
$\lambda_\pm(n_e)$ for SM-LMA versus LMA-Dark, showing (per its own caption)
"an anti-crossing visible inside the solar-core band" for SM-LMA and "no
anti-crossing" for LMA-Dark ("LMA-Dark avoids the solar resonance in normal
matter"); a daytime $P_{ee}(E)$ comparison annotated with real Borexino
($pp$, $^7$Be, $pep$) and SNO ($^8$B) data points (used only as plot
annotations here, not fitted -- the fit itself is
\cref{sec:results-borexino-nsi}); and the $A_{DN}(E)$ comparison plot with
the Super-K band. The notebook's own stated conclusion: "Super-K's
$A_{DN}=+3.3\pm1.0\%>0$ disfavours LMA-Dark at $\sim3\sigma$," while
explicitly caveating that "NSI in the solar interior (daytime degeneracy)
is not yet implemented in \code{solar\_probability\_state}. The dark-angular
daytime curve is therefore not the true LMA-Dark prediction but its
first-octant approximation."

\section{\texttt{intrinsic2\_BSM\_extension\_NSI.ipynb}: cross-method validation}

\begin{implbox}
This notebook cross-validates the numerical (matrix-exponential) and
perturbative/analytical propagation paths against each other for an
NSI-active scenario (oscillation preset \code{"\_SM\_NUFIT52\_NO"} with
\code{NSI\_extension="nsi\_globalfit\_esteban2018"}), across Earth
propagation and atmosphere-to-surface propagation, with an absolute-error
acceptance threshold of $10^{-3}$.
\end{implbox}

\begin{resultbox}
\textbf{Confirmed, directly printed output} (only two of this notebook's
comparison cells have recorded numeric output in the saved file; see the
caveat box above):

\emph{Pure flavour state through the Earth}
($E\in[100,2\times10^4]$~MeV $\times$ 90 log-spaced points,
$\eta\in[0.05,1.50]$~rad $\times$ 90 points; \code{method="numerical"},
400 midpoint steps, vs.\ \code{method="analytical"}):
\begin{center}
\small
\begin{tabular}{@{}lr@{}}
\toprule
Quantity & Value \\
\midrule
max.\ absolute error & $2.770\times10^{-2}$ \\
mean absolute error & $7.026\times10^{-4}$ \\
max.\ relative error & $5.648\times10^{-1}$ \\
mean relative error & $5.429\times10^{-3}$ \\
comparison points & 71\,523 \\
absolute threshold & $1.000\times10^{-3}$ \\
\bottomrule
\end{tabular}
\end{center}

\emph{Atmosphere production to surface}
($E\in[10^2,10^5]$~MeV $\times$ 90, $\theta\in[0^\circ,180^\circ]\times$90;
\code{method="numerical"}, 600 steps, vs.\ \code{method="analytical"}, 6
cubic segments):
\begin{center}
\small
\begin{tabular}{@{}lr@{}}
\toprule
Quantity & Value \\
\midrule
max.\ absolute error & $1.651\times10^{-11}$ \\
mean absolute error & $2.198\times10^{-13}$ \\
max.\ relative error & $1.551\times10^{-9}$ \\
mean relative error & $1.616\times10^{-11}$ \\
comparison points & 43\,874 \\
absolute threshold & $1.000\times10^{-3}$ \\
\bottomrule
\end{tabular}
\end{center}
\end{resultbox}

\begin{notebox}
The Earth-crossing case's maximum \emph{relative} error, $0.5648$, looks
large in isolation, but the notebook's own methodology (its Section~0.3)
applies relative-error comparison only where the reference probability
exceeds $10^{-5}$, precisely to avoid exactly this kind of division-by-small-number
inflation; the maximum \emph{absolute} error, $2.77\times10^{-2}$, is the
physically meaningful figure and is the one the acceptance criterion is
actually built around. Both confirmed comparisons pass a reasonable reading
of the $10^{-3}$ absolute-error criterion (the atmosphere case by many
orders of magnitude, reflecting the atmosphere's much lower electron
density; the Earth case's mean absolute error, $7.0\times10^{-4}$, is below
threshold, while its \emph{maximum} absolute error, $2.77\times10^{-2}$,
exceeds it at the grid's most extreme, longest-baseline points -- consistent
with the first-order perturbative expansion's known degradation for long,
high-density segments rather than an NSI-specific defect).

The notebook's design also includes five further comparisons (SM-limit
closure, solar-surface \code{adiabatic\_exact} vs.\ numerical,
solar-to-detector, atmosphere-to-detector, and a final
\code{assert np.isfinite(...)} summary printing "All four physical
comparisons completed successfully.") whose cells have no recorded output
in the saved file, so their outcome cannot be confirmed from the notebook
artifact as saved; re-running the notebook would be needed to obtain them.
\end{notebox}

\section{\texttt{inference2\_borexino\_nsi.ipynb}: fitting the degeneracy to real data}
\label{sec:results-borexino-nsi}

\begin{implbox}
This is the one NSI notebook that fits real data: it performs a
gradient-based, joint $(\theta_{12},\epsilon_{ee})$ fit
(\pyclass{SolarNSIOscillationModel}, \cref{sec:impl-autograd}, via
\pyfunc{inference.fit.fit\_lbfgs}) to the four real Borexino $P_{ee}(E)$
points published in~\parencite{Esteban2018LMADark}-style analyses (loaded
from \pkgpath{data/solar/borexino/probability/nature2018\_pee\_points.csv}),
starting from two deliberately different initial guesses -- one on the
standard-LMA side ($\theta_{12}=25^\circ$, $\epsilon_{ee}=+0.5$) and one on
the LMA-Dark side ($\theta_{12}=65^\circ$, $\epsilon_{ee}=-1.2$) -- to test
whether both converge to the degenerate solutions predicted by
\S\ref{sec:theory-lma-dark}'s analytic transformation. The base oscillation
preset used is \code{"\_SM\_NUFIT61\_NO"}.
\end{implbox}

\begin{resultbox}
\textbf{Confirmed, directly printed output} -- the degeneracy cross-check
comparing the dark-side LBFGS fit against the value predicted by applying
$\theta_{12}\to90^\circ-\theta_{12}$,
$\epsilon_{ee}\to-(2+\epsilon_{ee})$ to the independently-converged
LMA-side fit:
\begin{center}
\small
\begin{tabular}{@{}lrrr@{}}
\toprule
Quantity & Predicted from LMA fit & Actual dark-side fit & Difference \\
\midrule
$\theta_{12}^{\rm dark}$ [deg] & $55.50$ & $55.50$ & $+0.00$ \\
$\epsilon_{ee}^{\rm dark}$ & $-1.775$ & $-1.782$ & $+0.007$ \\
\bottomrule
\end{tabular}
\end{center}
Both printed as within $1\sigma$ of the predicted value
(\code{theta12\_dark within 1 sigma of prediction: True},
\code{eps\_ee\_dark within 1 sigma of prediction: True}).

\textbf{Confirmed, directly printed output} -- the independent 2-D
$\Delta(-2\ln\mathcal L)$ grid scan (LMA-side branch,
$30\times30=900$ points, $\theta_{12}\in[5^\circ,85^\circ]$,
$\epsilon_{ee}\in[-3.0,1.0]$):
\begin{equation*}
  \chi^2_{\min} = 0.777 \quad\text{at}\quad \theta_{12}=35.34^\circ,\ \epsilon_{ee}=-0.24
  \qquad(4\text{ data points}, 2\text{ free parameters}).
\end{equation*}
\end{resultbox}

\begin{notebox}
The exact numerical values of the LBFGS fits themselves (Section~4 of the
notebook, including its own \code{assert delta\_chi2 < 1.0} and "PASSED"
message) have no recorded cell output in the saved file, so they are
\textbf{not independently confirmable} from the artifact as saved. They can,
however, be partially back-derived from the confirmed Section~5 table above:
since $\theta_{12}^{\rm dark,\,predicted}=90^\circ-\theta_{12}^{\rm LMA}=55.50^\circ$,
the LMA-side fit must have converged to $\theta_{12}^{\rm LMA}\approx34.50^\circ$;
since $\epsilon_{ee}^{\rm dark,\,predicted}=-(2+\epsilon_{ee}^{\rm LMA})=-1.775$,
$\epsilon_{ee}^{\rm LMA}\approx-0.225$. These derived values are reasonably
consistent with the independently confirmed grid-scan minimum,
$\theta_{12}=35.34^\circ$, $\epsilon_{ee}=-0.24$ (a coarse $30\times30$ grid
evaluated independently of the gradient fit), which is expected agreement
between two different optimisation methods rather than a coincidence, but
this reasoning is presented here as \emph{derived}, not as a number the
notebook itself printed.
\end{notebox}

The notebook produces three figures: a filled $\chi^2(\theta_{12},\epsilon_{ee})$
contour map over the full grid with the standard-LMA and LMA-Dark preset
points marked as reference stars; an overlay of the two best-fit $P_{ee}(E)$
curves against the real Borexino data points with a residual-difference
panel; and the $\Delta(-2\ln\mathcal L)$ uncertainty map with Wilks'
68\%/90\%/95\% contours (at $\Delta\chi^2=2.30/4.61/6.18$ for 2 degrees of
freedom) and the grid minimum marked. The notebook's own stated conclusion:
"two deliberately different, uninformative starting guesses converge to two
well-separated points...with statistically indistinguishable
$\chi^2$...quantitatively consistent with the
$\theta_{12}\to90^\circ-\theta_{12}$,
$\epsilon_{ee}\to-(2+\epsilon_{ee})$ relation," and "the two fitted
$P_{ee}(E)$ curves agree to a fraction of a percent everywhere from 0.15 to
16~MeV" -- but also, echoing \cref{sec:results-lma-dark}, that "a real
resolution of LMA vs.\ LMA-Dark needs independent information...not more
solar $P_{ee}(E)$ points alone" (i.e.\ the Earth day--night asymmetry of
\cref{sec:results-lma-dark}, or another observable outside the solar
$P_{ee}(E)$ degeneracy itself).

\section{The composition-dependence extension: no notebook coverage yet}

\begin{notebox}
\textbf{This is the most important caveat in this chapter.} An exhaustive
search of every notebook in the repository (grep for \code{eps\_ee\_n},
\code{eps\_.*\_n}, \code{epsilon\_n}, \code{has\_neutron\_coupling}, and
\code{composition}) confirms that \textbf{no notebook currently
demonstrates the neutron-density composition-dependence extension}
($\epsilon^n$, \cref{sec:theory-composition,sec:impl-hamiltonian}). The
handful of unrelated matches (an Earth/solar/AK135 raw-data column literally
named \code{neutron\_density\_mol\_cm3}, evolutor-composition consistency
checks unrelated to NSI, PMNS flavour-composition language, etc.) are all
false positives on the search terms, not partial coverage. This is
consistent with the extension's implementation history: it was added after
every existing notebook was last executed, so no notebook artifact can yet
be cited as numerical evidence for it. The automated test suite
(\pkgpath{core/BSM/test/test2\_bsm\_nsi.py},
\pkgpath{core/common/test/test3\_hamiltonian.py},
\pkgpath{core/perturbative/test/test2\_perturbative\_evolutor.py}, and the
\pkgpath{medium/earth}, \pkgpath{medium/atmosphere}, \pkgpath{medium/solar}
test suites) is, at the time of writing, the only executed and verified
evidence for the composition-dependence extension's correctness; a natural
follow-up to this document would be a fifth notebook demonstrating it end
to end -- e.g.\ propagating a solar neutrino from the hydrogen-rich envelope
($n_n/n_e\to0$) through the helium/metal core ($n_n/n_e\to1$) with a
non-zero $\epsilon^n$, and comparing against the $\epsilon^n=0$
approximation used throughout this chapter's existing notebooks.
\end{notebox}
